\section{The Propers: Detailed Commentary}
\zlabel{triptych:concise:commentary:start}

\subsection{What changes when a debtor is released?}
The king's first act of judgment reveals an impossibility. The servant owes ten thousand talents and has nothing with which to pay. His promise to repay everything is therefore not a plan whose success the king can reasonably await. When the king releases him and cancels the debt, compassion supplies what an extension of time would never supply. Chrysostom reads the threat of sale through this outcome: it brings the man to supplication and awareness of what he owes; the answer exceeds his request. The gift does not originate in a merciful action the servant has already performed.\footnote{Chrysostom, \work{Homilies on Matthew} LXI, \S\S1--3, NPNF I.10.}

That order establishes the reading centered on mercy. The servant encounters his neighbor as a recipient before he acts as a creditor. His own empty hands before the king should change how he regards the other man's empty hands. Chrysostom makes the immense first debt intelligible by recalling creation and redemption, the world given to humanity, the Son sent after human ingratitude, baptismal remission and the Spirit. These are not entries in a balance sheet on which grateful creatures eventually cancel God's claim. Even their subsequent service is a life made possible by his gifts.

The Introit's creator and the Collect's family supply this argument with its beginning. Mordecai addresses the Lord of all from within a threatened people. The words do not make the petitioner an owner who can dispose of God's power. The Psalm verse's undefiled walk likewise names a dependence as well as a duty. Augustine hears the later plea for directed steps within that same Psalm: a command to walk rightly and a prayer for help belong to one voice. The Collect's good works follow continuing mercy and protection. They do not replace the need for them.\footnote{Esth.~13:8--18; Augustine, Ps.~118, Aleph \S\S1--8, especially 3,5--7.}

Augustine's debtor-and-creditor account brings the parable into the present tense. A hearer may admit once having needed pardon and then imagine that this condition is over. Augustine describes a person who still begs from God while another begs from him. We ask for Christ, the living bread, while a neighbor may ask us for bread. Giving and forgiving express a shared dependence that cannot be divided into a religious relation with God and a self-sufficient relation with everyone else.\footnote{Augustine, Sermon~83 \S\S1--4, NPNF I.6, Sermon XXXIII.}

The pilgrimage reading begins at the same release and asks a different question. Where does the liberated person go? The Alleluia celebrates Israel's departure from Egypt, while Augustine follows Paul's account of the sea, spiritual food and drink, and the subsequent falls in the wilderness. Baptismal deliverance is real; it does not make the remainder of the journey unreal. The king's servant illustrates the danger from within: he leaves bondage and uses that freedom to bind another. The mercy reading follows the obligation created by the gift; the pilgrimage reading follows the unfinished life begun by it. They converge in charity without becoming one argument.\footnote{Augustine, Ps.~113 \S\S3--6, with 1 Cor.~10:1--6; Chrysostom, \work{Ephesians} XXIII.}

The repeated plea makes the servant's failure unmistakable. His companion falls down and asks for patience in the very words that had brought him before the king. Chrysostom asks why he does not recognize the means by which he himself escaped. The contrast is more exact than rich against poor: the king granted complete remission when asked for time; the first servant denies even time after receiving remission. The lesser debt remains real, and injury by a neighbor can remain serious. What the man refuses is recognition of a fellow servant whose need now stands before him.

\subsection{How can real pardon end in judgment?}
The king's reproach refers to a gift actually made: he forgave all the debt. Chrysostom says that the recipient himself revoked the sentence through his cruelty; Augustine warns the baptized that what was forgiven may be exacted again. Their language gives the ending its force. A reading which makes the first pardon only a pretense would destroy the contrast: there would have been no received mercy for the servant to contradict. A reading which dismisses the final judgment would destroy the warning addressed to the disciples.\footnote{Chrysostom, \work{Matthew} LXI \S4; Augustine, Sermon~83 \S\S6--7.}

Nor does ``until'' repayment promise an eventual escape. Chrysostom reasons from the liability already established: the servant cannot discharge it. The duration of the imprisonment is not a calendar date concealed in the story. Jesus' final demand for forgiveness from the heart carries the hearer back from imagined accounting to the act now possible. The servant could not pay the king; he could have shown compassion to his companion.

Aquinas later distinguishes what the parable's dramatic reversal might otherwise leave confused. A new human fault does not undo God's completed act of remission. The old sinful acts do not regain their former guilt as though absolution had never occurred. Yet a new mortal sin incurs a new debt of eternal punishment, and contempt of a greater gift makes that new fault more grievous. Real pardon and renewed accountability thus remain true together. Aquinas expressly rejects the idea that forgiveness was merely provisional and never actually effected by grace.\footnote{Aquinas, \work{Summa theologiae} III, q.~88, a.~1, answer and replies; English Dominican translation, second revised edition (1920).}

His further distinctions prevent a mechanical picture of the debt's return. Article~3 does not add all previous penalties to the new one as a numerical total. The new sin has its own gravity, aggravated by ingratitude for received favor. Article~2 explains why hatred of a brother is especially contrary to pardon: it opposes the very form of the favor bestowed, remission of another's debt. This is a later theological explanation of renewed liability, not the technical vocabulary of Chrysostom or Augustine. Their severe preaching continues to ask what a pardoned person is doing with mercy.\footnote{Aquinas, III, q.~88, aa.~2--4, especially a.~2, answer, and a.~3, answer and reply~3.}

The two whole-Mass readings draw different consequences. In the first, the threatened judgment exposes ingratitude: freedom received must become mercy toward another. In the second, it prevents protection from being mistaken for guaranteed perseverance. A person may have truly crossed the sea and still fall. The danger comes not only from a hostile world but from a will that consents to cruelty. Hope is therefore neither self-reliance nor a prediction that every recipient will certainly remain faithful.

Augustine guards the demand for inward pardon against an equally damaging conclusion. Forgiveness does not require abandoning correction. A rebuke may seek a person's good while the heart relinquishes hatred. His ancient disciplinary examples do not prescribe a modern response to every injury; the distinction does make room for truthful naming of wrong, restitution and protection. The Gospel condemns the servant's deliberate refusal. It does not equate the lingering distress of a wounded person with that refusal.\footnote{Augustine, Sermon~83 \S8; restitution and protective boundaries are applications of his distinction.}

Job now matters before any theory of suffering begins. He is already upright when the adversary strikes. To identify an afflicted neighbor with the condemned servant merely because both suffer would reverse the Offertory's opening declaration. The Mass places guilt and innocent suffering beside one another without making them interchangeable. Mercy toward an offender and care for the injured can therefore belong to the same moral life.

\subsection{Whom does the armor resist, and what does protection promise?}
Ephesians' military language follows a summons to renewed conduct: anger must not become sin, bitterness must yield to forgiveness, and believers must imitate Christ's love. Immediately before the appointed verses, masters are reminded that they and their servants have one heavenly Master. Chrysostom expressly recalls the unmerciful servant at that point. The armor does not interrupt the teaching about charity; it equips the community to live it.\footnote{Eph.~4:26--32;5:1--2;6:9; Chrysostom, \work{Ephesians} XXII, NPNF I.13.}

Its exclusion is decisive: the struggle is not against flesh and blood. In the mercy reading, revenge on a fellow debtor is a battle against the wrong opponent. In the pilgrimage reading, the spiritual conflict explains why a delivered people still needs help. Truth, justice, peace, faith, salvation and God's word are not ornaments awarded after danger ends. They are gifts for standing within it. Opposition to wrongdoing can be firm without seeking the wrongdoer's destruction as its final good.

Chrysostom's details resist reduction to a rigid inventory of symbols. The breastplate is the whole virtuous life. The footwear can mean readiness to serve the Gospel or readiness for departure. Faith first appears as a gift of power and then as confidence in promised goods against assault. On two expressions his historical editor disagrees: Chrysostom takes ``in the heavenly places'' as the heavenly goods at stake, whereas the editor defends a locative reading; his ``evil day'' is the present life, while the notes prefer a day of violent temptation and discuss a final outbreak. These are different explanations of the words even where they encourage the same vigilance.\footnote{Chrysostom, \work{Ephesians} XXII--XXIV; NPNF I.13, notes 465,468.}

The difference between confidence and immunity becomes sharpest in Job. The narrative first permits an attack on possessions, then on the body while sparing his life. The Offertory condenses the two scenes; its broad marginal reference to Job~1 does not remove the bodily trial's place in chapter~2. Gregory begins from Satan's request: the adversary wants God to put forth his hand because he has no autonomous power to strike. Evil intention belongs to the creature; permitted power remains under divine government. Permission does not make the intention good.\footnote{Job~1:1,12--19;2:3--7; Gregory, \work{Moralia} II.10.16--18.}

This is why the Introit can call God Lord of all without making Haman's murderous project righteous. The enemy is neither an equal god nor a puppet whose malice must be called divine virtue. Gregory follows the limits of the trial: God both lets the adversary act and reins him in. Job is in the enemy's hand in one respect and in his helper's hand in another. Bodily injury is real, but it does not establish that the helper has ceased to hold him.\footnote{Gregory, \work{Moralia} II.11.19; III.5.6--7.}

The phrase ``without cause'' tests this explanation. Gregory distinguishes an antecedent crime deserving punishment from a providential end through which virtue is manifested and increased. Job's affliction lacks the former, not every possible purpose. He also refuses to imagine Satan changing God's mind into anger: the language of being moved describes divine action through ordinary human speech, not a deity overcome by his creature's persuasion. The Gradual's eternal refuge gives the same trust its liturgical voice. Our circumstances change; the God to whom we turn does not pass through successive defeats and recoveries.\footnote{Gregory, \work{Moralia} III.3.3--4; Augustine, Ps.~89 \S\S2--4.}

This account does not disclose the purpose of a particular neighbor's bereavement or illness. Gregory reads a narrative which declares Job righteous and reveals the heavenly challenge. An observer of a present wound does not possess the same revelation about its cause. The reading of protected pilgrimage consequently sustains the plea for rescue. It neither asks sufferers to desire injury nor excuses withholding help. The mercy reading receives the same truth as a limit on accusation: pain is not a reliable measure of another person's guilt.

\subsection{Why do those who stand firm also need food?}
Chrysostom develops the journey from the girdle itself. A traveler gathers his clothing in readiness to leave; Israel ate the Passover girded, shod and carrying a staff. Christ is the Christian Passover, and the Christian is likewise called to eat as one ready for departure. The comparison becomes an account of the Church's life: the sea has been crossed, the wilderness remains, the Spirit gives guidance, nourishment exceeds manna, and Christ leads toward heavenly rest. Armor and sacramental food serve the same unfinished passage.\footnote{Chrysostom, \work{Ephesians} XXIII, on 6:14 and the following Passover exposition.}

Augustine reaches this sequence through Psalm~113 and Paul's warning about Israel. His contribution gives the typology its moral seriousness: the historical deliverance truly prefigures baptism, yet those who received the signs could fall. Chrysostom presses a related consequence upon the hearer's possessions and habits. A pilgrim who lives for wealth as though the camp were the permanent city misunderstands the destination. The gift of food asks for a manner of life; it is not a reason to settle permanently within the danger it helps one traverse.

The Sunday Secret supplies the saving initiative of this journey. Propitiation means restoration of peace with God through sacrifice; the prayer asks for that reconciliation and renewed salvation through offerings whose effects God himself has willed. The family does not negotiate a saving price with a reluctant deity. Schuster's reading of Job toward the innocent Christ deepens the Offertory's passage into sacrifice: Christ bears our sins, while Job's literal uprightness still forbids the friends' accusation. The allegorical fulfillment does not turn every human wound into an expiation of its victim's own wrongdoing.\footnote{Schuster, \work{The Sacramentary} III (1927), pp.~178,181.}

The Trinity Preface prevents a divided account of this saving action. The Father, Son and Holy Spirit are confessed as distinct Persons in one divine essence and equal majesty. Augustine explains that the Son's mission does not imply inferiority to the Father, nor the Spirit's mission inferiority to either. The giving and receiving of salvation in history do not establish competing divine powers. The people already join angelic praise while still engaged in the Epistle's conflict.\footnote{Augustine, \work{On the Trinity} III, preface \S3; doctrinal illumination of the Preface.}

The Postcommunion makes the relation between gift and conduct particularly close. What the mouth has received must be followed with a pure mind. The historical English says ``receive with pure souls''; the Latin \textit{sectemur} also carries the force of following or pursuing. Schuster draws the sacramental bond into mind, heart and will. Augustine's demand for heartfelt pardon thus finds a concrete liturgical consequence: the communicant's inward life must be shaped by the gift, not left outside its reach.\footnote{Cummiskey (1861), XXI Sunday after Pentecost, Postcommunion; Schuster, III, p.~181.}

Here the two readings' spiritual senses remain distinct. In the mercy reading, Christ's redeemed household receives life which must become charity toward a fellow debtor. In the pilgrimage reading, the Church's Exodus in Christ receives food for a destination beyond present conflict. Both are allegorical readings anchored in baptismal and sacramental reception. Their moral fruits differ in emphasis---inward pardon and steadfast passage---while their final hope joins accountability to immortality. Neither treats the Eucharist as merely an image for human resolve.

\subsection{What is the Church asking when she asks for judgment?}
The Communion brings the Gospel's warning into the voice of people who suffer injustice. They hope in God's word, ask when he will judge their persecutors, and beg for help. Augustine hears the voice of the whole Christ: saints before the Incarnation longed for his coming, and saints after the Ascension await his manifestation as judge. Bellarmine places the emphasis on desire for everlasting security rather than curiosity about the world's remaining days. Personal petition opens toward the whole body's hope without making the present persecution disappear.\footnote{Augustine, Ps.~118, Caph, on vv.~81--84; Bellarmine, \work{Commentary on the Psalms}, Ps.~118:81--88.}

Augustine's explanation of the full Psalm's fainting helps a wounded hearer understand this hope. The soul fails toward salvation in longing, not away from it in abandonment. The appointed chant omits that opening verb but keeps salvation, hope and the cry for help. Emotional exhaustion is not therefore a sufficient sign that hope has vanished. The direction of desire matters: the person still reaches toward the God whose word he trusts.

Augustine then asks the hardest question of help: was prayer useless if martyrs were killed? He follows the next verses, in which the persecuted keep God's commandments and ask to be quickened by mercy. Death has not made them forsake the truth. Their life is preserved at the point where denial would have lost it; they live through dying for the truth. This answer gives the protection reading its final horizon. Rescue remains a real petition, but the salvation sought reaches beyond bodily survival.\footnote{Augustine, Ps.~118, Caph, on vv.~86--88.}

For the mercy reading, the same cry shows why forgiveness and judgment can coexist. The communicants entrust judgment to God; they do not strangle their persecutors. An end to injustice is a good to ask for. Relinquishing hatred does not require wishing persecution to continue, calling it harmless, or withdrawing protection from its victims. The Gospel judges mercilessness; it does not forbid the Communion's prayer.

Schuster puts a present judgment beside Augustine's and Bellarmine's final expectation. Rejection of truth already deprives its enemies of light. That is not simply another phrase for the final peace they await. Present privation explains a consequence already borne; final judgment answers the still-unfinished cry for justice and rest. The first reading therefore faces an ultimate account of mercy, while the second hopes for an end beyond the reach of persecution. The distinction makes their common hope more precise.\footnote{Schuster, III, p.~181.}

The dated prayers give that hope a communal and apostolic reach. In the mission Collect the gift goes outward through workers who speak confidently; in the Secret Christ's self-giving for all opens worship from sunrise to sunset; in the Postcommunion those already nourished ask that true faith continue to advance. The nations are recipients of proclamation, not the spiritual enemies against whom the armor is directed. Reception, mission and continuing fidelity form one movement of grace.

Where Luke is commemorated, the three prayers join a cross-bearing body to free service, healing, glory, sanctification and safety. The Collect's language is illuminated by 2 Corinthians~4, where outward affliction accompanies inward renewal and ministry points to Christ. It does not establish how Luke died. Heard together, the prayers keep safety from meaning absence of every wound and freedom from meaning independence of grace. A person fed at the altar can serve with liberty of mind while still carrying the cross. The family receives mercy to become merciful and strength to complete its journey; its end is lasting communion with the triune God.
